\documentclass[11pt,a4paper]{article}

\usepackage[margin=2.6cm]{geometry}
\usepackage{amsmath,amssymb}
\usepackage{booktabs}
\usepackage{graphicx}
\usepackage{xcolor}
\usepackage{enumitem}
\usepackage[colorlinks=true,linkcolor=blue!60!black,citecolor=blue!60!black,urlcolor=blue!60!black]{hyperref}

% Result placeholders: replace once the PointMaze run completes.
\newcommand{\TBD}[1]{\textcolor{red!70!black}{\textbf{[TBD: #1]}}}
\newcommand{\zlat}{\mathbf{z}}
\newcommand{\xloc}{\mathbf{x}}
\newcommand{\act}{\mathbf{a}}
\newcommand{\abar}{\bar{\mathbf{a}}}

\title{Addendum --- Local/Global World-Model Planning with DINO-WM:\\
A Decoupled Forward--Backward Planning Track\\
{\large Addendum to \emph{DINO-WM on PointMaze under a Commodity-GPU Budget};
implementation and protocol, results pending}}
\author{Thomas Georges\thanks{Implementation assistance: Claude (Anthropic). Code:
\texttt{World\_Models\_LAS} repository, \texttt{src/wm\_poc/local\_global}. This is an
addendum to the main report \emph{DINO-WM on PointMaze under a Commodity-GPU
Budget: Latent-Cached Reproduction, Low-Data Transfer, and Planning Results}
(\texttt{reports/dino\_wm\_report.tex}), whose trained checkpoint and latent
cache this track reuses unchanged.}}
\date{June 2026 --- addendum, draft v0.1}

\begin{document}
\maketitle

\noindent\emph{This is an addendum to the main report ``DINO-WM on PointMaze
under a Commodity-GPU Budget'' (\texttt{reports/dino\_wm\_report.tex}). That
report documents the DINO-WM model, its latent-cached training, and the
planning results this track builds on; its DINO-WM background and notation are
shared here and not re-derived.}
\vspace{0.6em}

\begin{abstract}
Accurate world models and differentiable world models pull in different
directions: the models that produce the most trustworthy forward rollouts
(large latent transformers, diffusion world models) are expensive or
impractical to differentiate through, while small differentiable surrogates
are easy to optimize against but easy to exploit. Following the decoupled
forward--backward idea of coupled local/global world models, we add a
\emph{local/global planning track} to an existing DINO-WM reproduction: the
frozen DINO-WM latent world model serves as the trusted \emph{global} forward
simulator and scorer, a small residual surrogate trained on cached DINOv2
patch latents serves as the \emph{local} differentiable model, and a family of
planners spans the spectrum between zero-order global search (CEM) and
first-order local optimization (gradient descent through the surrogate),
including hybrid planners in which global search proposes, local gradients
refine, and a global re-score accepts or rejects the refinement. We describe
the method, the offline latent goal-reaching evaluation protocol with its
fairness accounting, the safeguards against surrogate exploitation, and the
engineering required to run the full, un-watered-down experiment on freely
available hardware (Colab T4) with resumable training and evaluation.
Quantitative PointMaze results will be added in a subsequent revision.
\end{abstract}

\section{Introduction}
\label{sec:intro}

Model-predictive control with learned world models faces a structural
trade-off. Zero-order planners such as the cross-entropy method (CEM) only
require forward rollouts and therefore work with arbitrarily large and
accurate models, but their sample complexity grows quickly with horizon and
action dimensionality. First-order planners obtain gradients of the planning
objective with respect to actions and can be far more sample-efficient per
iteration, but they require a differentiable model --- and when that model is
small enough to differentiate cheaply, it is usually also inaccurate enough to
be \emph{exploited}: gradient descent happily walks into regions where the
surrogate is wrong.

The coupled local/global world-model line of work~\cite{fog2026} resolves the
trade-off by decoupling the two roles. A large, accurate \emph{global} model
produces forward trajectories and is never differentiated; a lightweight
\emph{local} model supplies derivatives, evaluated along the global model's
trajectories, so that the local model only ever needs single-step accuracy in
the neighborhood of states the global model visits. We transplant this idea
into a simpler, fully observable setting --- visual goal reaching in latent
space with DINO-WM~\cite{dinowm2025} --- where it becomes a statement about
\emph{planners} rather than policy learning:

\begin{itemize}[leftmargin=1.4em]
  \item the \textbf{global model} is the repository's existing DINO-WM latent
        world model (a frozen ViT predictor over frozen DINOv2 patch
        features~\cite{dinov2}), used as forward simulator and scorer, always
        under \texttt{no\_grad};
  \item the \textbf{local model} is a small residual MLP (or GRU) surrogate
        operating on a compressed projection of the same latents, trained from
        cached transitions, fully differentiable with respect to actions;
  \item \textbf{planners} combine the two: global CEM, local first-order
        (GD/Adam), local CEM (an ablation isolating model error from optimizer
        effects), and hybrid CEM~+~local refinement with an optional global
        re-score that can reject harmful refinements.
\end{itemize}

\paragraph{Contributions.} This report documents (i) a concrete
local/global planning instantiation on top of DINO-WM, including the surrogate
architecture and training objective (\S\ref{sec:method}); (ii) an offline
latent goal-reaching evaluation protocol with explicit fairness accounting and
honestly stated biases (\S\ref{sec:protocol}); (iii) three safeguards against
surrogate exploitation and a measurement of local--global disagreement
(\S\ref{sec:safeguards}); and (iv) the engineering that makes the full
experiment run, resume, and complete on commodity hardware
(\S\ref{sec:implementation}). Results tables and figures are left as
placeholders (\S\ref{sec:results}) to be filled once the first full PointMaze
run completes.

\section{Background}
\label{sec:background}

\paragraph{DINO-WM.} DINO-WM~\cite{dinowm2025} trains an action-conditioned
latent transition model directly on frozen DINOv2 patch features
$\zlat_t \in \mathbb{R}^{P\times D}$ ($P{=}196$ patches, $D{=}384$ for
ViT-S/14 at $224$\,px), with teacher-forced latent-consistency training and no
pixel reconstruction in the planning loop. Planning is visual goal reaching:
given a current and a goal observation, CEM/MPC minimizes the latent distance
between the predicted future state and the goal's encoding. The main report
(\emph{DINO-WM on PointMaze under a Commodity-GPU Budget},
\texttt{reports/dino\_wm\_report.tex}) reproduces DINO-WM training on PointMaze
with a precomputed latent cache; this addendum reuses that trained checkpoint
and cache unchanged.

\paragraph{Decoupled forward/backward models.} The FOG line of
work~\cite{fog2026} trains policies by differentiating short-horizon returns,
substituting the Jacobians of a lightweight RSSM for the true dynamics
derivatives while the trajectory itself is generated by a large frozen
diffusion world model. Two findings motivate our design: the small model alone
is insufficient as a forward simulator (their ``no diffusion'' ablation
fails), and the small model's derivatives are most useful when evaluated
\emph{at states produced by the trusted model}. Our MPC loop mirrors both: the
global model is the only simulator, and the surrogate re-anchors on globally
imagined latents at every replanning step.

\paragraph{Related planning paradigms.} PlaNet~\cite{planet} and
Dreamer~\cite{dreamer} plan or learn policies inside a learned recurrent
latent model; TD-MPC2~\cite{tdmpc2} combines a learned latent model with
sampling-based MPC and a value prior; differentiable MPC~\cite{diffmpc}
differentiates through the optimizer itself. V-JEPA~2~\cite{vjepa2} and the
JEPA world-model family~\cite{jepawm} support the broader pattern of planning
on top of frozen self-supervised representations that this track inherits.

\section{Problem setting and notation}
\label{sec:setting}

Episodes are sequences of latents $\zlat_0,\dots,\zlat_{T-1}$
($\zlat_t \in \mathbb{R}^{P\times D}$, cached fp16) with raw actions
$\act_0,\dots \in \mathbb{R}^{A}$ ($A{=}2$ for PointMaze, bounds
$[-1,1]^A$). Following DINO-WM, one \emph{model step} advances $f$ raw steps
(frameskip $f{=}5$): the model-step action is the folded block
\begin{equation}
\abar_k \;=\; \big[\act_{t_0+kf};\,\act_{t_0+kf+1};\dots;\act_{t_0+(k+1)f-1}\big]
\in \mathbb{R}^{fA},
\end{equation}
and latent frames are sampled at $t_0, t_0{+}f, t_0{+}2f,\dots$ Training
windows consist of $C$ context frames and $K$ rollout target frames
($C{=}2$, $K{=}3$), with the $C{-}1{+}K$ action blocks between consecutive
frames; a window starting at $t_0$ is valid iff
$t_0 \le \min\!\big(T_z - 1 - (C{+}K{-}1)f,\; T_a - (C{+}K{-}1)f\big)$.
Train/validation splits are by episode (never by window), drawn from a seeded
permutation of the \emph{uncapped} episode count so that membership is
invariant to smoke-test subsampling.

A \emph{goal-reaching task} is a tuple (episode, $t_0$): the planner observes
the context window ending at $t_{\mathrm{cur}}$ and must reach the latent
$\zlat_{\mathrm{goal}} = \zlat_{t_{\mathrm{cur}} + G f}$ recorded $G$ model
steps later in the same episode ($G{=}5$).

\section{Method}
\label{sec:method}

\subsection{Global forward model}
\label{sec:global}

The global model is the trained DINO-WM checkpoint, loaded in-process exactly
as the upstream planner loads it (Hydra train config + epoch checkpoint) and
driven on cached latents through the repository's latent-bypass patch, so
rollouts stay in latent space end to end. Three details matter:

\begin{enumerate}[leftmargin=1.4em]
  \item \textbf{Replay from an observed anchor.} The upstream rollout API
        expects one action block per context frame followed by future blocks.
        The adapter keeps the original \emph{observed} context (latents, raw
        proprioception, and the $C{-}1$ observed action blocks) plus the list
        of executed blocks, and replays from that anchor at every call. The
        model is therefore never conditioned on fabricated proprioception for
        imagined frames.
  \item \textbf{Normalization at the boundary.} Checkpoints trained with
        upstream action/proprio normalization receive normalized inputs inside
        the adapter (statistics recomputed from the raw dataset exactly as
        upstream computes them); every other component works in raw units.
  \item \textbf{Chunked candidate scoring.} CEM populations are scored in
        chunks of \texttt{rollout\_batch\_size} candidates. Candidates are
        independent, so chunking provably does not change scores or
        per-candidate call counts --- only peak memory --- which is what allows
        the identical experiment to run on a 16\,GB T4 (chunk 32) and an A100
        (chunk 128).
\end{enumerate}
All global-model calls execute under \texttt{torch.no\_grad()}; gradients can
only originate from the surrogate. For CPU-only tests and smokes, a synthetic
point-mass task with \emph{exactly} linear latents ($\zlat = W\mathbf{s} + b$)
provides a closed-form perfect global model (decode with the pseudo-inverse,
step the true dynamics, re-encode), so planner correctness is measured against
true dynamics rather than a learned proxy.

\subsection{Local surrogate}
\label{sec:local}

\paragraph{Projection.} Global latents are compressed by
$\xloc = W_{\!p}\,\mathrm{pool}(\mathrm{LN}(\zlat)) \in \mathbb{R}^{d}$
($d{=}256$), where pool is either a global patch mean or a $4{\times}4$
spatial grid average (preserving coarse object layout), and $W_{\!p}$ is a
\emph{frozen}, seeded, semi-orthogonal matrix. Freezing is principled rather
than incidental: the training loss is an MSE \emph{in projected space} between
predictions and projected targets, for which a jointly trained projector
admits the degenerate optimum $W_{\!p}{=}0$ (representation collapse). A
frozen near-orthogonal projection of LayerNormed features preserves distances
in the Johnson--Lindenstrauss sense and removes the collapse mode entirely; a
trainable variant exists behind a flag and requires an accompanying hinge
variance penalty.

\paragraph{Dynamics.} The default surrogate is a residual MLP,
$\xloc_{k+1} = \xloc_k + g_\theta([\xloc_k;\abar_k])$ (LayerNorm input, SiLU,
hidden width 512, 3 layers), so the identity map is the zero-function default.
A GRU variant consumes the $C{-}1$ context transitions with a single
\texttt{GRUCell} to build a hidden state (velocity inference) and predicts
residual deltas from it. Rollouts unroll the step function under autograd;
horizons are short ($\le 6$), so exact backpropagation through time is cheap.

\subsection{Surrogate training objective}
\label{sec:loss}

With $\hat{\xloc}_{1:K}$ the surrogate rollout from the context and
$\xloc_{1:K}$ the projections of cached target latents,
\begin{align}
\mathcal{L} \;=\;& \lambda_{\mathrm{roll}}\,
  \frac{\sum_{k=1}^{K} \gamma^{k-1}\,\lVert \hat{\xloc}_k - \xloc_k \rVert_2^2 / d}
       {\sum_{k=1}^{K} \gamma^{k-1}}
  \;+\; \lambda_{1}\,\lVert \hat{\xloc}_1 - \xloc_1 \rVert_2^2 / d \nonumber\\
  &+\; \lambda_{\Delta}\,\lVert (\hat{\xloc}_1 - \xloc_0) - (\xloc_1 - \xloc_0) \rVert_2^2 / d
  \;+\; \lambda_{J}\,\big\lVert \partial \hat{\xloc}_1 / \partial \abar_0 \big\rVert_2^2
  \;+\; \lambda_{V}\, \textstyle\sum_j \max(0, 1 - \sigma_j(\xloc)) ,
\label{eq:loss}
\end{align}
i.e.\ a (optionally discounted) multi-step rollout MSE, a one-step term, a
\emph{delta} term that keeps predicted state changes calibrated when absolute
error is dominated by static features, an optional action-Jacobian penalty
(double backward) reserved for taming exploitable gradients, and the variance
hinge used only with a trainable projector. Defaults:
$\lambda_{\mathrm{roll}}{=}\lambda_1{=}1$, $\lambda_\Delta{=}0.1$,
$\lambda_J{=}\lambda_V{=}0$, $\gamma{=}1$. Training reports a scale-free
diagnostic, the rollout MSE divided by the ``predict no change'' baseline;
values well below 1 certify the surrogate beats the trivial predictor
independently of the projection's numeric scale.

\subsection{Planners}
\label{sec:planners}

All planners optimize an open-loop sequence
$\abar_{1:H} \in \mathbb{R}^{H \times fA}$ toward
$\zlat_{\mathrm{goal}}$, minimizing
\begin{equation}
J(\abar_{1:H}) \;=\;
\underbrace{\tfrac{1}{PD}\lVert \hat{\zlat}_H(\abar_{1:H}) - \zlat_{\mathrm{goal}} \rVert_2^2}_{\text{goal cost}}
\;+\; \lambda_{\mathrm{sm}}
\underbrace{\tfrac{1}{(H-1)fA}\textstyle\sum_{k} \lVert \abar_{k+1}-\abar_k \rVert_2^2}_{\text{smoothness}} ,
\label{eq:objective}
\end{equation}
where $\hat{\zlat}_H$ is the model's predicted final latent (global model for
global planners; for local planners the analogous cost is computed in
projected space against $\xloc_{\mathrm{goal}}$). Planners optimize
$J$ but \emph{report} goal and smoothness components separately, so the
smoothness weight cannot contaminate cross-planner distance comparisons.

\paragraph{Global CEM.} Iterations sample a population of $N$ sequences from a
diagonal Gaussian (initial std expressed in units of half the action range),
clamp to bounds, score with the global model, and refit mean and std to the
top-$E$ elites with a std floor; the incumbent mean is always kept in the
pool, and the best-ever candidate (not the final mean) is returned. Sampling
uses a CPU generator seeded per planning round, making results
device-independent and reproducible.

\paragraph{Local first-order (GD/Adam).} Actions are parameterized through a
differentiable squash,
$\abar = \ell + (u-\ell)\,\tfrac{1}{2}(\tanh(\mathbf{r})+1)$, so iterates are
feasible by construction and gradients respect the boundary geometry. The raw
parameters are optimized by SGD or Adam with gradient-norm clipping, and the
best iterate (not the last) is returned. Warm starts from CEM solutions
require care: CEM clamps samples \emph{onto} the bounds, and
$\mathrm{atanh}(1-10^{-5})$ sits where the tanh derivative is
${\approx}2{\times}10^{-5}$, freezing refinement. Warm starts therefore
unsquash with a dedicated margin $\varepsilon_{\mathrm{ws}}{=}10^{-2}$
(derivative ${\approx}0.02$), trading a $\le 1\%$ bound offset for usable
gradients.

\paragraph{Hybrid CEM + local refinement (+ global re-score).} Global CEM
proposes $\abar^{\mathrm{CEM}}$ with global total cost $J^{\mathrm{CEM}}$;
the local planner refines from that warm start to $\abar^{\mathrm{ref}}$;
the global model re-scores the refinement \emph{on the same objective},
yielding $J^{\mathrm{ref}}$. With re-scoring enabled, the refinement is
rejected iff
\begin{equation}
J^{\mathrm{ref}} \;>\; J^{\mathrm{CEM}} + \tau\,\max(|J^{\mathrm{CEM}}|, \epsilon),
\qquad \tau = 0.05,
\label{eq:reject}
\end{equation}
in which case the CEM sequence is returned unchanged. Both
``refinement improved the global cost'' and ``refinement accepted'' are logged
per planning round. A \emph{local CEM} planner (CEM scored by the surrogate)
completes the $2{\times}2$ design (model $\times$ optimizer), isolating the
``small model'' effect from the ``gradients'' effect.

\paragraph{Compute accounting.} Every planner counts global forward rollouts
\emph{per candidate} (a CEM round of population $N$ over $I$ iterations costs
$NI{+}1$ global calls, the $+1$ being the final pure-goal re-evaluation),
local forward calls, backward steps, and wall time. These are first-class
columns of the results table: the local/global thesis is an
\emph{efficiency} claim, and the accounting is the currency it is paid in.

\subsection{Safeguards against surrogate exploitation}
\label{sec:safeguards}

Three safeguards are active by default (per the implementation
specification), with the spec's optional extensions deferred:
(1)~\emph{differentiable action bounds} via tanh squashing --- the first-order
planner cannot leave the feasible set, eliminating the most common
exploitation channel; (2)~an \emph{action smoothness penalty} in the
optimized objective; (3)~\emph{global re-scoring with rejection}
(Eq.~\ref{eq:reject}) --- the surrogate is never trusted to overrule the
global model. The evaluation additionally measures the surrogate's open-loop
\emph{disagreement}: the executed action sequence is replayed through the
surrogate and compared, step by step in projected space, against the global
model's imagined trajectory.

\section{Evaluation protocol}
\label{sec:protocol}

\paragraph{Offline latent goal reaching.} For each task, an MPC loop plans
$H' = \min(H, \text{steps remaining})$ blocks (the horizon shrinks toward the
goal), executes $m{=}1$ block on the \emph{global model as simulator}, appends
the imagined latent to the planning context, and repeats for $G$ steps.
Final performance is the global-latent distance to the goal. Per episode we
record:
\begin{itemize}[leftmargin=1.4em]
  \item \textbf{normalized final distance}
        $\;\rho = \lVert \zlat_{\mathrm{fin}} - \zlat_{\mathrm{goal}}\rVert^2 /
                  \lVert \zlat_{\mathrm{start}} - \zlat_{\mathrm{goal}}\rVert^2$,
        and \textbf{success} $\,= \mathbb{1}[\rho < 0.5]$ (``closed at least
        half the gap''), making thresholds scale-free across encoders;
  \item \textbf{reference replay}: the dataset's true action block sequence
        from start to goal replayed through the same simulator --- an
        achievable-by-construction calibration line (near zero under the exact
        synthetic model; its magnitude under DINO-WM measures the global
        model's own rollout error);
  \item local-space final distance, the disagreement metric of
        \S\ref{sec:safeguards}, refinement acceptance rate, bound-violation
        count (structurally zero), the three compute counters, and wall time.
\end{itemize}

\paragraph{Stated bias.} Because the global model is simulator \emph{and}
judge, \texttt{global\_cem} is structurally favored on distance metrics: it
optimizes exactly the quantity it is scored on. The protocol is therefore
read as follows: distance/success columns establish that local and hybrid
planners are \emph{competent}; the efficiency columns (wall time, call
counts) and the acceptance rate carry the comparative claim. Confirmation in
the real environment (MuJoCo PointMaze, via the host repository's upstream
planning path) is explicitly listed as the next step and is out of scope for
this revision.

\section{Implementation}
\label{sec:implementation}

The track follows the host repository's two-notebook convention: an
operational notebook (03) for setup verification, transition export,
surrogate training, and planner evaluation --- every heavy step delegated to a
gated CLI script --- and a read-only results notebook (08) that aggregates
artifacts into tables and figures. Engineering features relevant to
reproducing the experiment:

\begin{itemize}[leftmargin=1.4em]
  \item \textbf{Single experiment definition, hardware-only profiles.} The
        A100 config file defines the experiment (model sizes, optimizer
        budgets, episode counts, run name); the T4 config extends it and
        overrides \emph{throughput knobs only} (rollout chunk size, dataloader
        workers), inheriting the run name so a run started on one GPU tier
        resumes on another.
  \item \textbf{Resumability everywhere.} Surrogate training checkpoints
        model, optimizer state, and step counter every 1000 steps (plus
        best-by-validation every 500) onto persistent storage, and resumes by
        default; planning evaluation is wall-clock capped and resumes at
        per-episode granularity (the $i$-th task is deterministic in the
        seed), so the 100-episode sweep completes across sessions.
        An interruption costs at most 1000 training steps or one episode.
  \item \textbf{Determinism.} Episode splits, task sampling, CEM noise, and
        the frozen projection are all seeded; CEM sampling uses a CPU
        generator so results are identical across devices; chunked scoring is
        exactness-tested.
  \item \textbf{Testing.} ${\sim}170$ unit/integration tests run with torch
        (and a ${\sim}130$-test subset without it), including analytic toy
        planner tests --- notably a sign-flipped surrogate that makes local
        refinement provably harmful, asserting the rejection gate fires ---
        plus an end-to-end CPU smoke on the synthetic task covering all six
        planners in under two minutes.
\end{itemize}

\section{Experimental setup}
\label{sec:setup}

\paragraph{Task and data.} PointMaze with the host repository's DINO-WM
setup: frozen DINOv2 ViT-S/14 at $224$\,px ($P{=}196$, $D{=}384$), frameskip
$f{=}5$, raw action bounds $[-1,1]^2$; cached latents for ${\sim}2200$
episodes; episode-level validation split fraction $0.1$, seed $42$.

\begin{table}[h]
\centering
\small
\resizebox{\textwidth}{!}{%
\begin{tabular}{@{}llll@{}}
\toprule
\textbf{Component} & \textbf{Setting} & \textbf{Component} & \textbf{Setting} \\
\midrule
Projection & grid pool $4{\times}4$, frozen, $d{=}256$ & CEM population / elites / iters & $300$ / $30$ / $10$ \\
Dynamics & residual MLP, width 512, 3 layers & CEM init.\ std & $0.5\times$ half-range \\
Context / rollout & $C{=}2$, $K{=}3$ & GD/Adam iters, lr & $100$, $0.05$ \\
Batch / steps / lr & $512$ / $20\,000$ / $3{\times}10^{-4}$ & Horizon / goal steps / exec & $5$ / $5$ / $1$ \\
Weight decay & $10^{-4}$ (AdamW) & Smoothness $\lambda_{\mathrm{sm}}$, tolerance $\tau$ & $0.01$, $0.05$ \\
Loss weights & $\lambda_{\mathrm{roll}}{=}\lambda_1{=}1$, $\lambda_\Delta{=}0.1$ & Episodes / success threshold & $100$ / $\rho{<}0.5$ \\
\bottomrule
\end{tabular}}
\caption{Full-experiment configuration (\texttt{pointmaze\_surrogate\_a100.yaml};
the T4 profile differs only in rollout chunk size and dataloader workers).}
\label{tab:config}
\end{table}

\paragraph{Budget.} Surrogate training is data-bound (each step moves
${\sim}375$\,MB of cached latents), estimated at $40$--$80$ min on A100 and
$1$--$2.5$\,h on T4 with a $300$-minute per-session cap and cross-session
resume; the six-planner evaluation is dominated by global CEM rollouts and is
similarly capped and resumable. The episode count is set to $100$ rather than
$50$ so the binomial success rate has a $95\%$ CI half-width of
${\approx}0.10$ rather than ${\approx}0.14$ (a compromise short of $200$'s
${\approx}0.07$ at twice the runtime); the per-episode resume carries the
sweep across sessions. Planners evaluated:
\texttt{global\_cem}, \texttt{local\_cem}, \texttt{local\_gd},
\texttt{local\_adam}, \texttt{hybrid\_cem\_local\_refine},
\texttt{hybrid\_cem\_local\_refine\_global\_rescore}.

\section{Results}
\label{sec:results}

\TBD{This section will be completed after the first full PointMaze run
(surrogate training to $20\,000$ steps + six-planner evaluation, 100 episodes
each).}

\subsection{Surrogate training}
\label{sec:results-surrogate}

\TBD{training/validation loss curves; per-horizon validation rollout MSE
(steps $1\ldots K$); final \emph{vs-static} ratio. Figures exported by
notebook 08 to \texttt{\_summary/figures/}.}

\begin{table}[h]
\centering
\small
\begin{tabular}{@{}lccc@{}}
\toprule
 & one-step MSE & $K$-step rollout MSE & rollout MSE / static baseline \\
\midrule
final validation & \TBD{} & \TBD{} & \TBD{} \\
\bottomrule
\end{tabular}
\caption{Surrogate validation summary (placeholder).}
\label{tab:surrogate}
\end{table}

\subsection{Planner comparison}
\label{sec:results-planners}

\begin{table}[h]
\centering
\scriptsize
\setlength{\tabcolsep}{4pt}
\renewcommand{\TBD}[1]{\textcolor{red!70!black}{$\cdot$}}% compact placeholder inside this table
\resizebox{\textwidth}{!}{%
\begin{tabular}{@{}lcccccccc@{}}
\toprule
Planner & succ. & norm.\ dist.\ $\rho$ & wall time/ep.\ (s) &
glob.\ fwd & loc.\ fwd & bwd steps & accept.\ rate & disagreement \\
\midrule
\texttt{global\_cem}                          & \TBD{} & \TBD{} & \TBD{} & \TBD{} & --     & --     & --     & --     \\
\texttt{local\_cem}                           & \TBD{} & \TBD{} & \TBD{} & --     & \TBD{} & --     & --     & \TBD{} \\
\texttt{local\_gd}                            & \TBD{} & \TBD{} & \TBD{} & --     & \TBD{} & \TBD{} & --     & \TBD{} \\
\texttt{local\_adam}                          & \TBD{} & \TBD{} & \TBD{} & --     & \TBD{} & \TBD{} & --     & \TBD{} \\
\texttt{hybrid\_cem\_local\_refine}           & \TBD{} & \TBD{} & \TBD{} & \TBD{} & \TBD{} & \TBD{} & \TBD{} & \TBD{} \\
\texttt{hybrid\_\ldots\_global\_rescore}      & \TBD{} & \TBD{} & \TBD{} & \TBD{} & \TBD{} & \TBD{} & \TBD{} & \TBD{} \\
\midrule
reference replay (dataset actions)            & --     & \TBD{} & --     & --     & --     & --     & --     & --     \\
\bottomrule
\end{tabular}}
\caption{Planner comparison on offline latent goal reaching, 100 held-out
episodes (placeholder). Distances are judged by the global model (see the
stated bias in \S\ref{sec:protocol}); efficiency columns carry the
local/global claim.}
\label{tab:planners}
\end{table}

\TBD{optimization traces (CEM best cost per iteration; refinement stage);
hybrid refinement outcomes (improved / within-tolerance / rejected); analysis
of acceptance rate vs.\ surrogate disagreement.}

\subsection{Synthetic-task sanity check (preliminary)}
\label{sec:results-synthetic}

On the synthetic point-mass task --- where the global model is exact and the
surrogate is trained for only a few hundred steps --- the smoke-scale runs
(2--3 episodes; not statistics, a correctness signal) reproduce the expected
qualitative ordering: all six planners succeed; \texttt{local\_adam} attains
the lowest final distances (a well-trained surrogate plus gradients beats
sampling at equal budget); hybrids improve on pure CEM and their refinements
are accepted; plain \texttt{local\_gd} is weakest. In the adversarial unit
test where the surrogate's action gain is sign-flipped, refinement reliably
worsens the global cost and the re-score gate rejects it, returning the CEM
sequence unchanged.

\section{Discussion and limitations}
\label{sec:discussion}

\paragraph{What the protocol can and cannot conclude.} The offline protocol
measures planner behavior \emph{under the global model's belief about the
world}. It can establish that first-order refinement through a cheap
surrogate recovers (or fails to recover) global-CEM-quality solutions at a
fraction of the global-model compute, and how often the surrogate's gradient
directions agree with the global model. It cannot, by construction, rank
planners against reality: that requires the real-environment evaluation
listed below. The reference-replay line partially de-confounds this by
exposing the global model's own rollout error on demonstrated behavior.

\paragraph{Surrogate capacity vs.\ exploitability.} The surrogate is
deliberately small and its projection frozen; the open question the results
must answer is whether PointMaze dynamics in pooled DINOv2 space are smooth
enough for useful gradients at this capacity. The Jacobian penalty
($\lambda_J$), the GRU context variant, and the grid-vs-mean pooling ablation
are the pre-built levers if acceptance rates come back low.

\paragraph{Single seed, single task.} The first run is \texttt{seed0} on
PointMaze only; seed replication and the PushT extension (blocked on a PushT
latent cache upstream) are future work by design.

\section{Future work}
\label{sec:future}

(1)~Real-environment confirmation of the planner ranking in MuJoCo PointMaze,
reusing the upstream planning path for the global baseline and adding an
environment-backed episode runner for local/hybrid planners; (2)~seed
replication and the built-in ablations (projection, context model, optimizer);
(3)~FOG-style online surrogate fine-tuning on global-model rollouts to keep
local gradients on-distribution under the current action distribution;
(4)~surrogate ensembles with disagreement penalties inside the planning cost;
(5)~PushT, where contact dynamics should stress the surrogate and the global
re-score safeguard is expected to matter most.

\section*{Reproducibility}

Code lives in the \texttt{World\_Models\_LAS} repository
(\texttt{src/wm\_poc/local\_global}, \texttt{scripts/local\_global},
\texttt{configs/local\_global}); notebook 03 runs the full pipeline on Colab
(T4 or A100) with self-gating, resumable stages; \texttt{run\_smoke.sh}
reproduces the synthetic end-to-end check on CPU in under two minutes. Design
rationale is documented in \texttt{docs/local\_global\_techniques.md}, the
operational guide in \texttt{docs/local\_global\_runbook\_next\_steps.md}, and
the original implementation brief in
\texttt{LOCAL\_GLOBAL\_DINO\_WM\_IMPLEMENTATION\_SPEC.md}.

\begin{thebibliography}{9}
\setlength{\itemsep}{2pt}

\bibitem{fog2026}
P.~Amigo, R.~Khorrambakht, N.~Mansard, L.~Righetti.
\emph{Coupled Local and Global World Models for Efficient First-Order
Model-Based Reinforcement Learning}. arXiv:2602.06219, 2026.

\bibitem{dinowm2025}
G.~Zhou et al.
\emph{DINO-WM: World Models on Pre-trained Visual Features Enable Zero-shot
Planning}. 2025.

\bibitem{dinov2}
M.~Oquab et al.
\emph{DINOv2: Learning Robust Visual Features without Supervision}.
arXiv:2304.07193, 2023.

\bibitem{planet}
D.~Hafner et al.
\emph{Learning Latent Dynamics for Planning from Pixels}. ICML, 2019.

\bibitem{dreamer}
D.~Hafner et al.
\emph{Dream to Control: Learning Behaviors by Latent Imagination}. ICLR, 2020.

\bibitem{daydreamer}
P.~Wu et al.
\emph{DayDreamer: World Models for Physical Robot Learning}. CoRL, 2023.

\bibitem{tdmpc2}
N.~Hansen, H.~Su, X.~Wang.
\emph{TD-MPC2: Scalable, Robust World Models for Continuous Control}.
ICLR, 2024.

\bibitem{diffmpc}
B.~Amos et al.
\emph{Differentiable MPC for End-to-end Planning and Control}. NeurIPS, 2018.

\bibitem{vjepa2}
M.~Assran et al.
\emph{V-JEPA 2: Self-Supervised Video Models Enable Understanding, Prediction
and Planning}. 2025.

\bibitem{jepawm}
\emph{JEPA World Models for Physical Planning}. 2026.
(Full citation to be completed.)

\end{thebibliography}

\end{document}
